% ===========================================================================
\section{信息只会在处理中损失}\label{sec:information-properties}
\warmth{0}\quad\whisper{先条件化，再展开，最后丢掉一个非负项。}

\begin{strand}
互信息可以看作共享的证据。它非负且对称，可以按变量逐个展开；对观测结果
做后处理，不可能凭空制造关于信源的新信息。
\end{strand}

\block{核心恒等式}
\begin{proposition}[互信息的性质]\label{prop:mi-properties}
对离散随机变量 $X,Y$，
\[
  \MI(X;Y)\ge0.
\]
等号成立当且仅当 $X$ 与 $Y$ 独立。此外，
\begin{align*}
  \MI(X;Y)&=\MI(Y;X),\\
  \MI(X;Y)&=\Ent(X)-\Ent(X\mid Y)
           =\Ent(Y)-\Ent(Y\mid X).
\end{align*}
\end{proposition}

\begin{proof}
由定义~\ref{def:mi}，
\[
  \MI(X;Y)=\KL(p_{XY}\Vert p_Xp_Y).
\]
信息不等式给出非负性；等号成立
恰好对应 $p_{XY}=p_Xp_Y$，也就是独立性。交换 $X,Y$ 不改变定义中的
对数比，因此互信息对称。最后两个熵恒等式来自命题~\ref{prop:chain}
中的链式法则。
\end{proof}

\begin{definition}[条件互信息]\label{def:conditional-mi}
已知 $Z$ 后，$X$ 与 $Y$ 的共享信息定义为
\[
  \MI(X;Y\mid Z)
  :=\sum_zp_Z(z)\KL\!\left(
    p_{XY\mid Z}(\,\cdot,\cdot\mid z)
    \,\middle\Vert\,
    p_{X\mid Z}(\,\cdot\mid z)p_{Y\mid Z}(\,\cdot\mid z)\right).
\]
因此 $\MI(X;Y\mid Z)\ge0$。
\end{definition}

\begin{proposition}[互信息的链式法则]\label{prop:mi-chain}
\[
  \MI(X_1,\ldots,X_n;Y)
  =\sum_{i=1}^n
    \MI(X_i;Y\mid X_1,\ldots,X_{i-1}).
\]
\end{proposition}

\begin{proof}
逐项使用条件互信息的熵表达式
\[
  \MI(U;Y\mid V)=\Ent(Y\mid V)-\Ent(Y\mid U,V),
\]
其中 $U=X_i$、$V=(X_1,\ldots,X_{i-1})$。于是求和中的相邻条件熵
发生望远镜消去：
\begin{align*}
  \sum_{i=1}^n\MI(X_i;Y\mid X_1,\ldots,X_{i-1})
  &=\sum_{i=1}^n\bigl[
    \Ent(Y\mid X_1,\ldots,X_{i-1})
    -\Ent(Y\mid X_1,\ldots,X_i)\bigr]\\
  &=\Ent(Y)-\Ent(Y\mid X_1,\ldots,X_n).
\end{align*}
最后一个表达式正是 $\MI(X_1,\ldots,X_n;Y)$。
\end{proof}

\begin{yourturn}
展开两个变量的情形：
\[
  \MI(X_1,X_2;Y)
  =\MI(X_1;Y)+
   \answerin[3.9cm]{\MI(X_2;Y\mid X_1)}.
\]
\recall{缺少的一项是 $\MI(X_2;Y\mid X_1)$。}
\end{yourturn}

\block{数据处理}

\begin{definition}[Markov 链]
若给定 $Y$ 后 $X$ 与 $Z$ 条件独立，就记作 $X\to Y\to Z$。等价地，
\[
  p_{Z\mid XY}(z\mid x,y)=p_{Z\mid Y}(z\mid y).
\]
\end{definition}

\begin{theorem}[数据处理不等式]\label{thm:dpi}
若 $X\to Y\to Z$ 构成 Markov 链，则
\[
  \MI(X;Z)\le\MI(X;Y).
\]
\end{theorem}

\begin{proof}
按两种顺序展开同一个互信息：
\[
  \MI(X;Y,Z)
  =\MI(X;Y)+\MI(X;Z\mid Y)
  =\MI(X;Z)+\MI(X;Y\mid Z).
\]
Markov 条件给出 $\MI(X;Z\mid Y)=0$，而条件互信息非负。因此
\[
  \MI(X;Y)
  =\MI(X;Z)+\MI(X;Y\mid Z)
  \ge\MI(X;Z).
\]
\end{proof}

\begin{corollary}[确定性处理]
对任意函数 $f$，
\[
  \MI(X;f(Y))\le\MI(X;Y).
\]
\end{corollary}

\begin{proof}
因为 $f(Y)$ 完全由 $Y$ 决定，所以 $X\to Y\to f(Y)$ 构成 Markov 链。
应用定理~\ref{thm:dpi} 即得结论。
\end{proof}

\begin{yourturn}
说出数据处理不等式证明中的两个关键动作：
\workspace[1]
\answerprose{（1）按两种顺序展开 $\MI(X;Y,Z)$；（2）用 Markov 条件
$\MI(X;Z\mid Y)=0$，再用非负性 $\MI(X;Y\mid Z)\ge0$。}
\recall{两个动作：（1）按两种顺序展开 $\MI(X;Y,Z)$；（2）使用
$\MI(X;Z\mid Y)=0$ 与 $\MI(X;Y\mid Z)\ge0$。}
\end{yourturn}
